%%
%% This is file `sample-sigconf.tex',
%% generated with the docstrip utility.
%%
%% The original source files were:
%%
%% samples.dtx  (with options: `sigconf')
%%
%% IMPORTANT NOTICE:
%%
%% For the copyright see the source file.
%%
%% Any modified versions of this file must be renamed
%% with new filenames distinct from sample-sigconf.tex.
%%
%% For distribution of the original source see the terms
%% for copying and modification in the file samples.dtx.
%%
%% This generated file may be distributed as long as the
%% original source files, as listed above, are part of the
%% same distribution. (The sources need not necessarily be
%% in the same archive or directory.)
%%
%%
%% Commands for TeXCount
%TC:macro \cite [option:text,text]
%TC:macro \citep [option:text,text]
%TC:macro \citet [option:text,text]
%TC:envir table 0 1
%TC:envir table* 0 1
%TC:envir tabular [ignore] word
%TC:envir displaymath 0 word
%TC:envir math 0 word
%TC:envir comment 0 0
%%
%%
%% The first command in your LaTeX source must be the \documentclass
%% command.
%%
%% For submission and review of your manuscript please change the
%% command to \documentclass[manuscript, screen, review]{acmart}.
%%
%% When submitting camera ready or to TAPS, please change the command
%% to \documentclass[sigconf]{acmart} or whichever template is required
%% for your publication.
%%

\documentclass[acmsmall,screen, nonacm, anonymous]{acmart}


\newcommand*{\ONECOLUMN}{}

% !TEX root = ../main.tex

\usepackage{docmute}
\usepackage{hyperref}
\usepackage{proof}

\usepackage{mathtools}
\usepackage{tikz}
\usepackage{tikz-network}
\usepackage{xspace}
\usepackage{xargs}
\usepackage{fontenc}

%\usepackage{minted}

\usepackage{subcaption}
\usepackage{caption}

\usepackage{tikz-cd}
\usepackage{tikzit}
% to reduce tikz picture white space
\usepackage{adjustbox}

% for quotienting
\usepackage{xfrac} 

\usepackage{import}

% pseudocode
\usepackage{algorithm}
\usepackage{algpseudocode}

% category names
%\newcommand{\MdaCospans}{{\catname{MACsp_{D}(Hyp_{\Sigma})}}}

\newcommand\HypI[1]{\textbf{HypI}(#1)}
\newcommand\MdaCospans{\textbf{MHypI}(\Sigma)}
\newcommand{\Ecospans}{{\catname{EHypI(\Sigma)}}}
\newcommand{\MdaEcospans}{{\catname{MEHypI(\Sigma)}}}
\newcommand{\WellTypedMdaEcospans}{{\catname{MEHypI(\Sigma)}}}
% conflict
\newcommand{\hashtag}{{\#}}
\newcommand{\consistency}{{\smile}}
% There is already a remark in the preamble, but it does
% not work for some reason
\newtheorem{remark}{Remark}


\usepackage{xifthen}
\ifthenelse{\isundefined{\ismain}}{
  \newcommand{\ismain}{0}
}{}

\usepackage{enumitem}
 
\newcommand{\ignora}[1]{ }

\ignora{
\usepackage{aliascnt}
% alias-counters
\newaliascnt{definition}{thm}
\newaliascnt{proposition}{thm} 
\newaliascnt{lemma}{thm}
\newaliascnt{corollary}{thm}
\newaliascnt{conjecture}{thm}
\newaliascnt{remark}{thm}
\newaliascnt{example}{thm}
\newaliascnt{examples}{thm}
\newaliascnt{assumption}{thm}
\newaliascnt{construction}{thm}
\newaliascnt{claim}{thm}

% \theoremstyle{plain}
\theoremstyle{definition}

\newtheorem{theorem}[thm]{Theorem}
\newtheorem{proposition}[proposition]{Proposition} 
\aliascntresetthe{proposition}
\newtheorem{lemma}[lemma]{Lemma}
\aliascntresetthe{lemma}
\newtheorem{corollary}[corollary]{Corollary}
\aliascntresetthe{corollary}
\newtheorem{conjecture}[conjecture]{Conjecture}
\aliascntresetthe{conjecture}
\newtheorem{remark}[remark]{Remark}
\aliascntresetthe{remark}
\newtheorem{assumption}[assumption]{Assumption}
\aliascntresetthe{assumption}
\newtheorem{construction}[construction]{Construction}
\aliascntresetthe{construction}
\newtheorem{claim}[claim]{Claim}
\aliascntresetthe{claim}

\theoremstyle{definition}
\newtheorem{definition}[definition]{Definition}
\aliascntresetthe{definition}
\newtheorem{example}[example]{Example}
\aliascntresetthe{example}
\newtheorem{examples}[examples]{Examples}
\aliascntresetthe{examples}
}

\usepackage{cleveref} % after hyperref!
% possibly new thm environment
% \newtheorem{satz}[thm]{Satz}


% \newcommand{\TODO}[1]{{{\color{red}[\textbf{TODO:} #1]}}}
% \newcommand{\TODO}[1]{}

% !TEX root = ../main.tex

\newcommand{\missing}[1]{{\color{red}\bfseries [TODO]}}

\newcommand{\egg}{\texorpdfstring{\MakeLowercase{\texttt{egg}}}{\texttt{egg}}\xspace}
\newcommand{\Egg}{\texorpdfstring{\MakeLowercase{\texttt{egg}}}{\texttt{egg}}\xspace}
% \newcommand{\egg}{Lego\xspace}
% \newcommand{\Egg}{Lego\xspace}
\newcommand{\egraphs}{\mbox{e-graphs}\xspace}
\newcommand{\egraph}{\mbox{e-graph}\xspace}
\newcommand{\Egraph}{\mbox{E-graph}\xspace}
\newcommand{\Egraphs}{\mbox{E-graphs}\xspace}
\newcommand{\eclass}{\mbox{e-class}\xspace}
\newcommand{\Eclass}{\mbox{E-class}\xspace}
\newcommand{\enode}{\mbox{e-node}\xspace}
\newcommand{\eclasses}{\mbox{e-classes}\xspace}
\newcommand{\enodes}{\mbox{e-nodes}\xspace}
\newcommand{\Enodes}{\mbox{E-nodes}\xspace}
% \newcommand{\regraph}{Regraph\xspace}
% \newcommand{\Regraph}{Regraph\xspace}
\newcommand{\sz}{Szalinski\xspace}
\newcommand{\find}{\texttt{find}\xspace}

\newcommand{\equivid}{\equiv_{\sf id}}
\newcommand{\equivnode}{\equiv_{\sf node}}
\newcommand{\equivterm}{\equiv_{\sf term}}

\newcommand{\congrinv}{$\mathcal{I}_c$\xspace}

\newcommand{\egglogo}[1][]{\protect\includegraphics[height=1em, #1]{egg.png} }
\newcommand{\eggurl}{\url{https://github.com/mwillsey/egg}}

% TODOs
% https://tex.stackexchange.com/questions/9796/how-to-add-todo-notes

\newcommand{\eggtodo}[1]{\textcolor{Red}{\textbf{[CITE: #1]}}}

\newcommand{\defTodo}[2]{%
  \expandafter\newcommand\csname #1\endcsname[1]{%
    \todo[linecolor=#2,backgroundcolor=#2!25,bordercolor=#2,inline,size=\tiny]{\textbf{#1}: ##1}}}

\newcommand{\defTODO}[2]{%
  \expandafter\newcommand\csname #1\endcsname[1]{%
    \todo[linecolor=#2,backgroundcolor=#2!25,bordercolor=#2,inline,size=\tiny,caption={\textbf{(#1 LONG TODO)}}]{##1}}}

\newcommand{\AlekseiColor}{RoyalBlue}

% examples
% \newcommand{\RemyColor}{Red}
% \newcommand{\OliverColor}{OliveGreen}
% \newcommand{\ChandraColor}{Magenta}
% \newcommand{\PavelColor}{Cerulean}
% \newcommand{\ZachColor}{Plum}
% \newcommand{\BenColor}{Orchid}
% \newcommand{\JamesColor}{Salmon}


\defTodo{Aleksei}{\AlekseiColor}
% \defTodo{Max}{\MaxColor}

\defTODO{ALEKSEI}{\AlekseiColor}
% \defTODO{MAX}{\MaxColor}


% categorical stuff

\newcommand{\catname}[1]{\mathbf{#1}}

\newcommand{\Graph}{\catname{Graph}}
\newcommand{\Set}{\catname{Set}}

\newcommand\id{\textsf{id}}
\newcommand\sym{\textsf{sym}}

% tiny inline string diagrams
\newcommand{\comonoid}{%
	\begin{tikzpicture}[baseline=0.4ex, scale=0.4, line width=0.8pt]
		\path [use as bounding box] (0,0) rectangle (0.6,0.8);
		\draw (0.2,0) -- (0.2,0.4);
		\filldraw[black] (0.2,0.4) circle (0.08);
		\draw (0,0.8) .. controls (0,0.65) and (0.1,0.45) .. (0.2,0.4);
		\draw (0.4,0.8) .. controls (0.4,0.65) and (0.3,0.45) .. (0.2,0.4);
	\end{tikzpicture}%
}
\newcommand{\counit}{%
	\begin{tikzpicture}[baseline=0.4ex, scale=0.4, line width=0.8pt]
		\path [use as bounding box] (0,0) rectangle (0.4,0.8);
		\draw (0.2,0) -- (0.2,0.4);
		\filldraw[black] (0.2,0.4) circle (0.08);
	\end{tikzpicture}%
}

% camera-ready
\newcommand{\ifcameraready}[2]{\ifdefined\cameraready #2 \else #1 \fi}

\input{../../sample.tikzstyles}
\input{../../hypergraph.tikzstyles}
\input{../../hypergraph.tikzdefs}


\definecolor{applegreen}{rgb}{0.55, 0.71, 0.0}
\definecolor{americanrose}{rgb}{1.0, 0.01, 0.24}
\definecolor{atomictangerine}{rgb}{1.0, 0.6, 0.4}
\definecolor{azure}{rgb}{0.0, 0.5, 1.0}

% These are for comments
\newcommand\question[1]{{\color{azure}#1}}
\newcommand\update[1]{{\color{americanrose}#1}}


\usepackage{minted}
\usepackage{csquotes}

% type derivations
\usepackage{mathpartir}


\makeatletter
\newsavebox{\@brx}
\newcommand{\llangle}[1][]{\savebox{\@brx}{\(\m@th{#1\langle}\)}%
  \mathopen{\copy\@brx\kern-0.5\wd\@brx\usebox{\@brx}}}
\newcommand{\rrangle}[1][]{\savebox{\@brx}{\(\m@th{#1\rangle}\)}%
  \mathclose{\copy\@brx\kern-0.5\wd\@brx\usebox{\@brx}}}
\makeatother


%%
%% \BibTeX command to typeset BibTeX logo in the docs
\AtBeginDocument{%
  \providecommand\BibTeX{{%
    Bib\TeX}}}


\settopmatter{printacmref=false}
\authorsaddresses{}

\begin{document}

\title{Notes on category theory}

\author{Anonymous}


\maketitle

\section{Yoneda Lemma}

The lemma~\cite{maclane} sometimes pops up in various proofs of category theory statements.

\begin{remark}

The following notation can be used for a set of natural transformations between two functors
\[
F,G : \mathbf{C} \to \mathbf{D}
\]

\[
\alpha : \Pi_{A : |\mathbf{C}|} \mathbf{D}(FA,GA)
\]
or,
\[
\int_{A : \catname{C}} \catname{D}(FA,GA)
\]

\end{remark}

\begin{lemma}[Yoneda]

Let $F : \catname{C}^{\text{op}} \to \catname{Set}$ be a contravariant functor to $\catname{Set}$ (i.e. a \textit{presheaf}).
Recalling that $\catname{C}(-, A)$ is a contravariant functor:
\begin{itemize}
\item $\catname{C}(-,A)B = \catname{C}(A,B)$
\item $\catname{C}(-,A)f : \catname{C}(D,A) \to \catname{C}(C,A) = g \circ f$ for $f : \catname{C}(C,D)$ and $g : \catname{C}(D,A)$
\end{itemize}

We can define a functor $Y : \catname{C} \to \catname{PShC}$ by letting
\begin{itemize}
\item $YA \coloneq \catname{C}(-,A)$
\item $Yf : \catname{C}(-,A) \to \catname{C}(-,B) \coloneq \int_{X : \catname{C}}\catname{C}(X,A) \to \catname{C}(X,B)$ for $f : \catname{C}(A,B)$, $(Yf)_{X}g = f \circ g$
\end{itemize}

Then we have for every contravariant functor $F : \catname{C}^{\text{op}} \to \catname{Set}$
\[
\int_{Y : \catname{C}}(YX)Y \to FY \cong FX
\]
or,
\[
\int_{Y : \catname{C}}\catname{C}(Y,X) \to FY \cong FX
\]
\end{lemma}

\begin{proof}

First let's define 
\[
\psi_{X} : (\int_{Y : \catname{C}}\catname{C}(Y,X) \to FY\;) \to FX
\]

as
\[
\psi_{X} \alpha \coloneq \alpha_{X} \circ id_{X}\;,
\]
i.e. each natural transformation $\alpha$ is mapped to $\alpha_X \circ id_{X}$ (into its component in $X$ applied to $id_{X}$)

and then
\[
\phi_{X} : FX \to (\int_{Y : \catname{C}}\catname{C}(Y,X) \to FY \;)
\]
as
\[
\phi_{X}\; x \coloneq \lambda f : Y \to X . (Ff x)\;.
\]
Let's check if these are the inverses of each other.
\begin{align*}
(\phi_{X};\psi_{X}) x &= \psi_{X}(\phi_{X} x)\\ 
 &= (F\;id_{X}) x\\
 &= id_{F_{X}} x\\
 &= x
\end{align*}

\begin{align*}
 (\psi_{X};\phi_{X}) \alpha &= \lambda f . (Ff) (\alpha_{X} \circ id_{X})\\
  &\overset{\text{naturality of } \alpha}{=} \lambda f . \alpha_{Y} \circ (\catname{C}(f,X) \circ id_{X})\\
  &= \lambda f . \alpha_{Y} \circ (id_{X} \circ f)\\
  &= \lambda f . \alpha_{Y} \circ f\\
  &= \alpha
\end{align*}
\end{proof}

\begin{corollary}
    The functor $Y : \catname{C} \to \catname{PShC}$ is full and faithful.
\end{corollary}
\begin{proof}

Using the lemma by plugging $Y$:
\[
\int_{Z : |\catname{C}|} \catname{C}(Z,X) \to YAZ \cong YAX
\]
or
\[
    \int_{Z : |\catname{C}|} \catname{C}(Z,X) \to \catname{C}(Z,A) \cong \catname{C}(X,A)
\]
We can note that the left-hand side is the action of $Y$ on morphisms of $\catname{C}$.
\end{proof}

To show that $f : \catname{C}(A,B)$ is an isomorphism we can check if $YA \cong YB$, \textit{i.e.}, if $\catname{C}(-,A) \cong \catname{C}(-,B)$ in a sense $\int_{X : |\catname{C}|} \catname{C}(X,A) \cong \catname{C}(X,B)$.

\begin{remark}
    The above is the contravariant version of the lemma.
\end{remark}

Recall the free-forgetful adjunction between $\catname{Mon}$ and $\catname{Set}$ where the free functor $F$ maps a set $X$ to a monoid $FX$ consisting of strings $x_1,x_2,\ldots,x_n$ and the monoidal multiplication is given by string concatenation.
Functor $U$ is then maps any given monoid $M$ to the set of its underlying elements $UM$.
We then have that for every set-morphism $X \to UM$ there is a unique monoid-morphism $FX \to M$.
\[
\catname{Mon}(FX, M) \cong \catname{Set}(X, UM)
\]
For example, let's take a monoid of natural numbers $\mathbb{N}$ and a set $X = \{a,b,c\}$.
We can define a function $f : \{a,b,c\} \to \mathbb{N}$ by $\{a \mapsto 1, b \mapsto 2, c \mapsto 3 \}$.
A free monoid over $\{a,b,c\}$ is $FX = \{\varepsilon,a,b,c,aa,bb,cc,\ldots\}$.
This induces a function $g : FX \to M$ by defining its action on generators $\{a \mapsto 1, b \mapsto 2, c \mapsto 3, \varepsilon \mapsto \varepsilon\}$ and then $g(aa) = g(a)\cdot g(a)$ and so on.

We can analogously define a similar adjunction between $\catname{Cat}$ and $\catname{Graph}$ (a directed graph is given by a set of vertices and edges such that for every edge we have a source and a target vertex).
Every category $\catname{C}$ defines a graph $U\catname{C}$ with the same objects and arrows forgetting which are composites and which are identities.
Let an $O$-graph be a graph over a set of vertices $O$.
A morphism of two $O$-graphs is a function that is identity on vertices.
Any $O$-graph can be used to generate a category $C$ on the same set of objects where the morphisms are all `strings' of composable arrows of the graph.

\begin{theorem}

Given a small graph $G = \{A \rightrightarrows O\}$.
There is a small category $\catname{C}$ with objects $O$ and a morphism $P : G \to U\catname{C}$ of graphs with the following property.
Given a category $\catname{B}$  and any morphism $D : G \to U\catname{B}$ there is a unique functor $D' : \catname{C} \to \catname{B}$ with $(U D') \circ P = D$.
In that sense $P$ is `universal' among the morphisms $G \to U\catname{C}$.
\end{theorem}
\begin{proof}

Let the objects of $\catname{C}$ be the ones of $\catname{G}$ and the arrows be the finite paths $a_1 \xrightarrow{f_1} a_2 \xrightarrow{} \ldots \xrightarrow{f_{n-1}} a_{n}$.
The desired morphism $P : G \to U\catname{C}$ is identity on objects and sends an arrow $a_1 \xrightarrow{f_{1}} a_2$ of $G$ to $\langle a_1, f_1, a_{2} \rangle$ of $U\catname{C}$.
That is, $\catname{C}$ is a free category over $G$.
Now, consider any other morphism $D : G \to U\catname{B}$ for a category $\catname{B}$, such that there is a functor $D' : \catname{C} \to \catname{B}$.
If $(UD') \circ P = D$, then $D_{O}(G) = D'_{O}(C)$ and $D'\langle a_{1}, f_{1}, a_{2} \rangle = Df_{1}$. 
\end{proof}
\begin{example}
    \item When a graph $G$ is given by a single vertex $a$ and a single arrow $a \xrightarrow{f} a$, the morphisms of $FG = \catname{C}$ are $1, f, f^2, \ldots$ (i.e. $\langle a \rangle, \langle a , f, a \rangle, \langle a f a f a \rangle, \ldots$)
\end{example}

\begin{definition}

If $S : \catname{D} \to \catname{C}$ is a functor and $c$ is an object of $\catname{C}$, a \textit{universal} arrow from $c$ to $S$ is a pair $\langle r, u \rangle$ consisting of an object $r$ of $\catname{D}$ and an arrow $u : c \to Sr$ of $\catname{C}$ such that to every pair $\langle d, f \rangle$ with $d$ an object of $\catname{D}$ and $f : c \to Sd$ there is a unique arrow $f' : r \to d$ such that $Sf' \circ u = f$.
That is, every arrow $f'$ is factorized uniquely through $u$.
\end{definition}
\begin{example}

Take $U : \catname{Cat} \to \catname{Graph}$ and let $c$ be an object of $\catname{Graph}$.
The universal morphism is then $\langle r, u \rangle$ where $r = Fc$ and $u : c \to U(Fc)$.
Then, given any other $\langle d, f \rangle$, where $f : c \to Ud$, we have that there exists $j : Fc \to d$ such that $f = u \circ Uj$
\end{example}

\begin{definition}[$T$-algebra]

Given a monad $\langle T : X \to X, \eta : I \to T, \mu : T \circ T \to T \rangle$, a $T$-algebra is given by a pair $\langle x, h \rangle$ where $h : Tx \to x$ and $x \in X$, such that
\[
    \begin{tikzcd}
        T^2x \arrow[r, "Th"] \arrow[d, "\mu_{x}"'] & Tx \arrow[d, "h"] & x \arrow[r, "\eta_{x}"] \arrow[rd, no head, Rightarrow] & Tx \arrow[d, "h"] \\
        Tx \arrow[r, "h"]                          & x                 &                                                         & x                
    \end{tikzcd}
\]
A morphism between two $T$-algebras $\langle x , h \rangle$ and $\langle y, h' \rangle$ is given by a morphism $f : x \to y$ such that $h' \circ Tf = f \circ h$.
\end{definition}

\begin{theorem}[Every monad is defined by its $T$-algebra]
If $\langle T, \mu, \eta \rangle$ is a monad on $X$, then the set of $T$-algebras and their morphisms forms a category $X^T$ and there is an adjunction
\[
\langle F^T, G^T, \eta^{T}, \varepsilon^{T} \rangle : X \to X^{T}
\], where the functors are given by
\[
    \begin{tikzcd}
        {\langle x, h \rangle} \arrow[r, "{G^{T}\langle x, h \rangle}"] \arrow[d, "f"] & x \arrow[d, "G^{T}f"] & x \arrow[r, "F^{T}x"] \arrow[d, "f"'] & {\langle Tx, \mu_{x} \rangle} \arrow[d, "F^{T}f"] \\
        {\langle y, h' \rangle} \arrow[r, "{G^{T}\langle y, h \rangle}"]               & y                     & y \arrow[r, "F^{T}y"]                 & {\langle Ty, \mu_{y} \rangle}                   
    \end{tikzcd}
\]
and $\eta^{T} = \eta : I \to G^{T} \circ F^{T}$ and $\varepsilon^{T}\langle x, h \rangle = h$.
Recall the typing $\varepsilon : F^{T} \circ G^{T} \to I$, or, component-wise $\varepsilon \;\langle F^T \circ G^T \rangle \langle x,h \rangle = \langle x,h \rangle$, $\varepsilon \;\langle Tx, \mu_{x}\rangle = \langle x,h \rangle$.
That is, the co-unit is the morphism between $T$-algebras $\langle Tx, \mu_{x} \rangle$ and $\langle x, h \rangle$.
That is, this is just the morphism between the carriers $Tx \to x$, i.e. $h$.
The monad defined by this adjunction is the given monad $T$.
First, let's check if the triangle equations hold
\[
    \begin{tikzcd}
        F \arrow[r, "F\eta"] \arrow[rd, no head, Rightarrow] & F \circ G \circ F \arrow[d, "\varepsilon_{F}"] &  & G \arrow[r, "\eta_{G}"] \arrow[rd, no head, Rightarrow] & G \circ F \circ G \arrow[d, "G\varepsilon"] \\
                                                             & F                                              &  &                                                         & G                                          
    \end{tikzcd}
\]
We get
\[
\begin{tikzcd}
        {Fx = \langle Tx, \mu_{x} \rangle} \arrow[rr, "F\eta"] \arrow[rd, no head, Rightarrow] &                               & {F(Tx) = \langle TTx, \mu_{Tx} \rangle} \arrow[ld, "\varepsilon_{F}"] &  & {G\langle x, h \rangle = x} \arrow[r, "\eta_{G}"] \arrow[rd, no head, Rightarrow] & Tx \arrow[d, "G_{\varepsilon}"] \\
                                                                                               & {\langle Tx, \mu_{x} \rangle} &                                                                       &  &                                                                                   & x                              
\end{tikzcd}
\]
The right-most arrow of the right-hand diagram is given by $\varepsilon_{\langle x, h \rangle} = h : Tx \to x$, i.e. the morphism $\langle Tx, \mu_{x} \rangle \to \langle x, h \rangle$ and $G\varepsilon : Tx \to x$.
The endofunctor $(G^{T} \circ F^{T})x = G^{T}(\langle Tx, \mu_{x} \rangle) = Tx$ defines the original monad $T$.
The unit of the monad is $\eta$ and the join $\mu\;x : TTx \to Tx$ is given by $G^{T}\varepsilon F^{T}\;x = G^{T}\varepsilon(\langle Tx, \mu_{x}\rangle) = G^{T}\mu_{x} = \mu_{x}$ ($G^{T}$ just `lifts' $\mu_{x}$ from homomorphism of $T$-algebras to a plain morphism in $X$).
\end{theorem}

\begin{theorem}[Comparison of adjunctions with algebras]

Given an adjunction $X \to A$, $\langle F, G, \eta, \epsilon \rangle$ and a monad on $X$ given by $\langle G \circ F, \eta, G\varepsilon F \rangle$.
Then there is a unique functor $K : A \to X^{T}$ such that $G^{T}K = G$ and $KF = F^{T}$, that is the following diagram commutes
\[
\begin{tikzcd}
        A \arrow[r, "K", dashed] \arrow[d, "G", shift left=2] & X^{T} \arrow[d, "G^{T}", shift left] \\
        X \arrow[u, "F"] \arrow[r, no head, Rightarrow]       & X \arrow[u, "F^{T}", shift left]    
\end{tikzcd}
\]
\end{theorem}
\begin{proof}
 The co-unit of the adjunction gives us the morphism $G\varepsilon_{a} : GFGa \to Ga$ and the arrow may be considered as the evaluation map of the corresponding $T$-algebra $\langle Ga, G\varepsilon_{a} \rangle$.
 This is a valid $T$-algebra by checking that the necessary diagrams commute.
 We define $Ka = \langle Ga, G\varepsilon_{a} \rangle$ and $Kf : \langle Ga, G\varepsilon_{a} \rangle \to \langle Gb, G\varepsilon_{b} \rangle$.
 Now we can check, \textit{e.g.}, that $KF = F^{T}$, $KFa = \langle GFa, G\varepsilon_{Fa} \rangle$, $F^{T}a = \langle GFa, \mu_{a}\rangle$. Recalling that $\mu_{a} = G\varepsilon_{Fa}$ we get the equality.
\end{proof}

\begin{definition}

When the functor $K$ is an isomorphism, we say that $G$ is \textit{monadic}.
\end{definition}

\begin{example}

Consider a free-forgetful adjunction $\langle F,G,\eta,\varepsilon \rangle : \catname{Set} \to \catname{SmGrp}$.
$FX$ is defined by turning a set $X$ into a free semigroup, i.e. $x \mapsto \langle x \rangle$ where the multiplication is given by words juxtaposition.
The forgetful functor simply forgets the multiplication and maps $\langle x \rangle \mapsto x$.
$\eta : X \to GFX$ sends each $x$ to one-letter word $x$.
The co-unit $\varepsilon : FGS \to S$, the composite $G\varepsilon \circ \eta_{G}(S) : GS \to GFGS \to GS$ is then the identity, in other words $\varepsilon$ maps generators.
For example, let $X = \{a,b,c\}$, $FX$ is then $\{\langle a \rangle, \langle b \rangle, \langle c \rangle, \langle aa \rangle, \ldots\}$.
$GFX = \{a,b,c,aa,\ldots\}$ and $\eta$ maps $a \mapsto a, b \mapsto b, c \mapsto c$.
Given $S = \{\langle s \rangle, \langle ss \rangle, \langle sss \rangle, \ldots\}$, $GS = \{s, ss, sss, \ldots\}$, $FGS = \{\langle s \rangle, \langle ss \rangle, \langle sss \rangle, \ldots\}$ and $\varepsilon$ just maps $\langle s \rangle \mapsto \langle s \rangle$ etc.
In other words $\varepsilon(\langle s_1 \rangle \ldots \langle s_n \rangle) = s_1 \cdot \ldots \cdot s_n$.

The induced monad $T = GF$ such that $TX = \coprod_{n=1}^{\inf}X^{n}$, $\eta_{x} = \langle x \rangle$, 
\[
\mu = G\varepsilon_{F}(\langle \langle x_{11} \rangle, \ldots, \langle x_{1,n_1} \rangle, \ldots, \langle \langle x_{k,1}\rangle, \ldots, \langle x_{k,n_{k}} \rangle \rangle) = \langle x_{1,1}\rangle \ldots \langle x_{k,n_{k}} \rangle~.
\]

For the above monad the algebras are $\langle S, v \rangle$ where $v : TS \to S$ i.e. $v : \coprod S^{n} \to S$ which is given by $\coprod_{1}^{\inf} v^{n}$.
The law $v\eta_{x} = id$ implies that $v_{1} = id$.
We also have that $v_{k}(v_{n_1} \times \ldots \times v_{n_k}) = v_{n_{1} + \ldots n_{k}} : S^{n_{1} + \ldots + n_{k}} \to S$, i.e. applying $k$ operations to a list of length $k$ that consists of lists of lengths $n_1$ to $n_k$ is the same as applying the operation to a list of length $n_1 + \ldots + n_{k}$.

The isomorphism $K : \catname{SmGrp} \to \catname{Set}^{T}$ is then given by $\langle S, v \rangle \mapsto \langle S, v_1, v_2, \ldots \rangle$ where $v_n$ is the iterated version of $v$.
\end{example}

\begin{example}

Consider now the adjunction $\langle F, G, \eta, \varepsilon \rangle : \catname{Graph} \to \catname{Cat}$.
The corresponding $T$-algebra for a monad $GF$ is $\langle Tx, \mu_{x}$ where $T$ adds identity paths and all other finite paths to a graph via designated edges.
The functor $K : \catname{Cat} \to \catname{Graph}^{T}$ maps a category $\catname{C}$ to $\langle G\catname{C}, h : GFG\catname{C} \to G\catname{C} = G\varepsilon_{C}\rangle$.
Given a category $\catname{C}$, $G\catname{C}$ is a graph with edges generated by morphisms of $\catname{C}$, then, $FG\catname{C}$ is a category where morphisms are all the finite paths of $G\catname{C}$ where each edge in the path is induced by a morphism in $\catname{C}$ and the whole path is given by the composition of morphisms, therefore we can map each path to a composition of the morphisms.
This defines the action of $\varepsilon$.
This is carried to the morphism of the respective underlying graphs and defines the structural map $h$.
Note that not any graph can be a component of a $T$-algebra.
In particular, it should be such graph $X$ that $TX \to X$ is definable, i.e. there must be an edge in $X$ for every finite path in $X$.
This means that there is a corresponding category for $X$ in $\catname{Cat}$ that has the same objects and morphisms are edges.
Since a graph has an edge for every path, it means that there is a composite morphism for any other two morphisms and the composition is associative.
\end{example}


\begin{definition}[Exponentials]
    In a category $\catname{C}$ with binary products we say there are \textit{exponentials} if for all objects $A$ $B$ there exist an object $A \multimap B$ and a map $ev_{A,B} : A \multimap B \times A \to B$ such that for any $g : C \times A \to B$ there is a unique morphism $\Lambda(g) : C \to (A \multimap B)$ such that the following diagram commutes
    \[
    \begin{tikzcd}
        A \multimap B \times A \arrow[r, "{ev_{A,B}}"]               & B \\
        C \times A \arrow[u, "\Lambda(g) \times id"] \arrow[ru, "g"] &  
    \end{tikzcd}
    \]
    We also have $\Lambda(ev_{A,B} \circ h \times id_{A}) = h$ for $h : C \to (A \multimap B)$.
    $h \times id_{A} : C \times A \to (A \multimap B) \times A$, $ev_{A,V} \circ h \times id_{A} : C \times A \to B$ and $\Lambda(ev_{A,B} \circ h \times id_{A}) : C \to A \multimap B = h$ by uniqueness.
\end{definition}

\bibliographystyle{ACM-Reference-Format}
\bibliography{../../bibliography}

\end{document}
